\documentclass[9pt,a4paper,onecolumn,twoside]{tau-class/tau}
\usepackage[english]{babel}
\usepackage{pgfplotstable}
\usepackage{booktabs}
%----------------------------------------------------------
% TITLE
%----------------------------------------------------------

\journalname{Example Template}
\title{Airbnb Pricing Factors in Columbus: A Statistical Analysis}

%----------------------------------------------------------
% AUTHORS, AFFILIATIONS AND PROFESSOR
%----------------------------------------------------------

\author[a,1]{Monish Sinha}

%----------------------------------------------------------

\affil[a]{STAT 325, Case Western Reserve University}

\professor{Professor Soale}

%----------------------------------------------------------
% FOOTER INFORMATION
%----------------------------------------------------------

\institution{Case Western Reserve University}
\footinfo{STAT 325 Project}
\theday{November 14, 2025}
\leadauthor{Sinha et al.}
\course{STAT 325}

%----------------------------------------------------------
% ABSTRACT AND KEYWORDS
%----------------------------------------------------------

\begin{abstract}    
This project explores the key factors that influence nightly Airbnb listing prices in Columbus, Ohio. Drawing on data from Inside Airbnb, this statistical analysis aims to help hosts optimize pricing while offering insights for policymakers about short-term rental impacts in a rapidly growing urban market.
\end{abstract}

%----------------------------------------------------------

\keywords{Airbnb, Columbus, regression, pricing, short-term rental, urban economics}

%----------------------------------------------------------

\begin{document}

\maketitle 
\thispagestyle{firststyle} 
\tauabstract 

%----------------------------------------------------------

%----------------------------------------------------------

\section{Introduction}

\taustart{T}he rapid growth of short-term rental platforms like Airbnb has fundamentally transformed urban housing markets and travel accommodation. In Columbus, Ohio---one of the fastest-growing cities in the Midwest---the Airbnb market has expanded significantly, raising questions about what factors drive nightly pricing and how hosts can optimize their listings. Understanding these pricing dynamics is valuable for multiple stakeholders: hosts seeking to maximize revenue, guests evaluating value, and policymakers assessing the economic impact of short-term rentals on local housing markets.

This study investigates the determinants of Airbnb nightly prices in Columbus using data from Inside Airbnb, a publicly available dataset containing detailed listing information. We focus on property characteristics (guest capacity, number of bedrooms), guest satisfaction metrics (review scores), and room type classification (entire home, private room, or shared room) as potential pricing factors. A key research question is whether different room types exhibit distinct pricing strategies---specifically, whether the relationship between guest capacity and price varies by room type.

We employ multiple linear regression with interaction terms to model nightly prices, comparing candidate models of increasing complexity to identify the specification that best balances explanatory power and parsimony. Our analysis reveals that room type significantly moderates the capacity--price relationship: private rooms show steeper per-guest price increases than entire homes, while shared rooms exhibit minimal capacity-based pricing. Additionally, guest review scores emerge as a strong predictor, with each additional rating point associated with substantial price premiums.

The remainder of this paper is organized as follows. Section~2 describes the data cleaning process and model specification. Section~3 presents the results, including model selection, coefficient estimates, and diagnostic checks. Section~4 discusses the implications of our findings and concludes with limitations and directions for future research.

\section{Methods}

\taustart{T}he dataset used for this project was downloaded from Inside Airbnb, specifically the detailed listings file for Columbus, Ohio (\texttt{listings.csv.gz}). This file contains a comprehensive snapshot of all active Airbnb listings in the city at the time of scraping and includes various fields such as pricing, room characteristics, location data, review counts, and host attributes.

The raw dataset includes mixed data types: numerical values stored as strings, date fields with inconsistent formatting, boolean values in 't' or 'f' format, and text-heavy columns like \texttt{amenities} encoded as JSON-like strings. Several fields also contained missing or malformed values, especially for price and location coordinates.

Data cleaning was performed using a custom R script \texttt{clean\_listings.R}, written to produce a cleaned version of the data in a reproducible and transparent way. There is a readme too with detailed instructions of the process. The cleaning process consisted of the following key steps:

\begin{itemize}
    \item Rows missing the unique listing ID were dropped, and duplicates were removed by retaining the first occurrence of each ID.
    \item Price values (e.g., \$128.00) were parsed and converted into numeric format under a new column, \texttt{price\_num}.
    \item Numeric columns such as \texttt{accommodates}, \texttt{bathrooms}, \texttt{bedrooms}, \texttt{beds}, and \texttt{number\_of\_reviews} were coerced into appropriate numeric formats.
    \item Boolean flags (e.g., \texttt{host\_is\_superhost}) were standardized from 't'/'f' strings into logical \texttt{TRUE}/\texttt{FALSE}/\texttt{NA}.
    \item Important date fields like \texttt{last\_scraped}, \texttt{host\_since}, and \texttt{first\_review} were parsed into Date objects with consistent formatting.
    \item Coordinates were normalized and stored as numeric values (\texttt{latitude}, \texttt{longitude}).
    \item The \texttt{amenities} column was lightly cleaned by removing excessive whitespace and newlines, producing \texttt{amenities\_clean}.
    \item Flags were created for rows with missing prices or coordinates to allow optional exclusion or further processing downstream.
    \item A full column-level report of data types and missing values was generated and saved to support  specification development.
\end{itemize}

The cleaned dataset was saved as \texttt{listings\_clean.csv}, retaining a subset of core variables needed for modeling. This ensures that downstream regression modeling is based on consistent, numeric-ready data. Rows with missing prices or invalid coordinates were retained with flags, allowing for flexible decisions on imputation or removal in future steps.

The data-cleaning script also automatically installs required packages (\texttt{readr}, \texttt{dplyr}, \texttt{lubridate}, \texttt{stringr}, \texttt{tibble}) and supports rerunning from the terminal or RStudio. The process is designed to be reproducible as possible if new columns or transformations are needed.

\subsection{Final Model Specification and Coefficient Interpretation}

The selected model specification was:
\[
\text{Price} = \beta_0 + \beta_1 \text{accommodates} + \beta_2 \text{bedrooms} + \beta_3 \text{review\_scores} + \beta_4 \text{review\_activity} + \text{room type effects}
\]

where room type effects include main effects for private and shared rooms (relative to entire homes) and their interactions with guest capacity.

All main effects were highly significant (p < 0.001). For each additional guest the property accommodates, nightly price increased by approximately \$14.28 for entire homes. Each additional bedroom added approximately \$15.42 to the nightly rate. Guest review scores demonstrated a strong price premium: a one-point increase on the (typically) 5-point satisfaction scale was associated with a \$31.68 increase in nightly price. Review activity (measured in reviews per month) carried a small negative coefficient (-\$1.46 per review/month), though this was marginally significant (p = 0.057), potentially reflecting a market effect where high-volume turnovers occur at lower price points.

The interaction term revealed the most interesting market dynamic: private rooms showed a much steeper relationship between guest capacity and price compared to entire homes. The interaction coefficient of $19.57 indicates that the price increase per additional guest is $19.57 higher for private rooms than for entire homes. Given the baseline accommodates effect of $14.28 for entire homes, this implies a per-guest slope of approximately $33.86 for private rooms, compared with $14.28 for entire homes. In contrast, the shared-room interaction term was small and non-significant, yielding an implied slope of only about $1.60 per guest. These results are consistent with Figure~\ref{fig:interaction}, where private rooms display the steepest price–capacity slope, entire homes a moderate slope, and shared rooms a nearly flat pattern.

The model explained approximately 48\% of the variance in nightly prices ($R^2 = 0.484$), with a root-mean-squared error of approximately \$71. This level of fit is reasonable given market complexity; unmeasured factors such as location amenities, host reputation, and dynamic market demand account for the remaining variation.

%----------------------------------------------------------

\section{Results}

\taustart{A}fter removing listings with missing prices or review scores, 2,381 listings remained for analysis. The nightly price ranged from \$23 to \$997, with a mean of \$135 and median of \$108, suggesting a right-skewed distribution typical of short-term rental markets. We focused our analysis on understanding how property characteristics, guest satisfaction, and room type influence pricing.

\subsection{Exploratory Model Development and Selection}

To identify the pricing mechanisms at work in Columbus's Airbnb market, we fit and compared four candidate linear regression models with increasing complexity. Our initial model (Model 1) included the main effects of guest capacity (accommodates), number of bedrooms, guest review satisfaction (review\_scores\_rating), and review frequency (reviews\_per\_month), as well as room type classification.

Before settling on this baseline model, we recognized a potentially important nuance in the market structure. In short-term rental markets, different room types---entire homes versus private rooms---often operate under different pricing strategies and may scale differently as capacity increases. To test this hypothesis, we developed Model 2 by adding an interaction between room type and guest capacity. The interaction term was motivated by the observation that entire homes, with more space to accommodate guests, might command steeper per-guest premiums compared to private rooms, which have inherent spatial constraints.

We also tested an interaction between room type and review scores (Model 3) to capture whether guest satisfaction drives price premiums differently across room types. Finally, Model 4 included multiple interactions to explore maximum explanatory power. Model comparison using the Akaike Information Criterion (AIC) and adjusted $R^2$ revealed that Model 2 provided the best balance between model fit and parsimony. Model 2 achieved an adjusted $R^2$ of 0.483 and AIC of 27,095, compared to Model 1's adjusted $R^2$ of 0.469 and AIC of 27,154. Model 4, while achieving higher $R^2$ (0.511), suffered from inflated AIC (26,976) and was more complex; Model 3 showed no improvement over the baseline.

Formal hypothesis testing via ANOVA confirmed that the room type by accommodates interaction in Model 2 was statistically significant (F(2,2372) = 31.65, p < 0.001), justifying its inclusion. In contrast, the room type by review scores interaction in Model 3 did not reach significance (p = 0.967), suggesting guest satisfaction affects pricing uniformly across room types.

\begin{table}[htbp]
\centering
\caption{ANOVA Variance Decomposition}
\label{tab:anova}
\begin{tabular}{lrrrrr}
\hline
Term & Df & Sum Sq & Mean Sq & F value & Pr(>F) \\
\hline
accommodates              &    1 & 10588790.41 & 10588790.41 & 2075.7895 & 0      \\
bedrooms                  &    1 &   205233.46 &   205233.46 &   40.2333 & 0      \\
review\_scores\_rating    &    1 &   220961.37 &   220961.37 &   43.3165 & 0      \\
review\_activity          &    1 &    20155.84 &    20155.84 &    3.9513 & 0.047  \\
room\_type                &    2 &     4417.70 &     2208.85 &    0.4330 & 0.6486 \\
accommodates:room\_type   &    2 &   322918.01 &   161459.01 &   31.6519 & 0      \\
Residuals                 & 2372 & 12099786.94 &     5101.09 &        -- & --     \\
\hline
\end{tabular}
\end{table}

\begin{table}[htbp]
\centering
\caption{Regression Parameter Estimates}
\label{tab:params}
\begin{tabular}{lrrrrc}
\hline
Parameter & Estimate & Std. Error & t value & Pr(>|t|) & Sig. \\
\hline
(Intercept)                               & -116.60855 & 21.226538  &  -5.493526 & 0        & *** \\
accommodates                              &  14.282121 &  0.952232  &  14.998575 & 0        & *** \\
bedrooms                                  &  15.420057 &  2.325319  &   6.631371 & 0        & *** \\
review\_scores\_rating                    &  31.678952 &  4.491330  &   7.053356 & 0        & *** \\
review\_activity                          &  -1.458741 &  0.767359  &  -1.900989 & 0.057425 &     \\
room\_typePrivate room                    & -39.496159 &  6.923357  &  -5.704770 & 0        & *** \\
room\_typeShared room                     & -10.842444 & 123.762666 &  -0.087607 & 0.930197 &     \\
accommodates:room\_typePrivate room       &  19.573072 &  2.460586  &   7.954639 & 0        & *** \\
accommodates:room\_typeShared room        & -12.677371 & 87.483494  &  -0.144912 & 0.884793 &     \\
\hline
\end{tabular}
\end{table}

\subsection{Model Assumptions and Diagnostics}

We assessed the linear regression assumptions to ensure the reliability of our coefficient estimates and hypothesis tests. Visual inspection of residuals versus fitted values revealed no clear patterns of nonlinearity, though some heteroscedasticity was evident---a common feature in price data where variability tends to increase with predicted price. A Breusch--Pagan test formally confirmed heteroscedasticity (BP = 283.64, p < 0.001), suggesting that prediction intervals around higher prices are wider than around lower prices. A Shapiro--Wilk normality test indicated slight departure from normality in the residuals (W = 0.809, p < 0.001), with some heavy right tails reflecting occasional extreme high-price outliers. However, with our large sample size (n = 2,381), these violations are unlikely to severely impact coefficient estimates or significance tests due to the Central Limit Theorem.

Multicollinearity was not a concern; all variance inflation factors were below 5 (the highest was 4.52 for accommodates), indicating that predictors were sufficiently independent. Influential outliers were identified via leverage plots but represented a small fraction of the data; robustness checks (not shown) confirmed that key conclusions remained stable even when excluding extreme observations.

\subsection{Visualization of Key Findings}

\begin{figure}[h]
\centering
\includegraphics[width=\linewidth]{figures/06_interaction_effect_plot.png}
\caption{Predicted nightly price as a function of guest capacity (accommodates) for each room type under Model 2. Private rooms show the steepest increase in predicted price as the number of guests increases, indicating that private-room listings in this dataset scale nightly prices more aggressively with capacity. Entire homes/apartments show a moderate upward trend—still increasing with capacity, but at a slower rate than private rooms. Shared rooms remain nearly flat with minimal price change as guest capacity increases. All predictions hold bedrooms, review scores, and review activity at their median values. This visualization reflects how the interaction between room type and guest capacity leads to different pricing behaviors across the three room types.}
\label{fig:interaction}
\end{figure}

The interaction effect is illustrated in Figure~\ref{fig:interaction}, which displays predicted prices across guest capacities for each room type. The diverging slopes confirm that private rooms exhibit a steeper price–capacity relationship than entire homes, with predicted prices increasing more rapidly as capacity rises. Entire homes show a more moderate positive slope, while shared rooms remain nearly flat across guest capacities. This pattern highlights how different room types scale their pricing strategies with capacity in distinct ways.

\subsection{Model Diagnostics and Supplementary Figures}

To validate the reliability of our regression inferences, we examined standard diagnostic measures and residual behavior during model development. Pairwise relationships among predictors and the response were reviewed to ensure no unexpected patterns or strong nonlinearities. These exploratory checks confirmed that the bivariate associations aligned with the structure of our multivariate model.

Although standard diagnostic plots were examined during model development, they were removed for brevity. Residual patterns indicated some heteroscedasticity and mild departures from normality, primarily driven by a small number of high-price listings. These issues are common in pricing data and, given the large sample size, are unlikely to meaningfully affect coefficient estimates or inference. Overall, the model meets the linear regression assumptions sufficiently well for the purposes of this analysis.

%----------------------------------------------------------

\section{Discussion and Conclusion}

\subsection{Key Findings and Interpretation}

\taustart{T}he room type interaction reveals sophisticated market segmentation. Entire homes serve different market segments (families, groups) than private rooms (single travelers, couples), and the pricing structure reflects these different use cases. Private rooms’ steeper scaling with guest capacity suggests that, in this dataset, hosts raise prices more aggressively as they host additional guests, whereas entire homes show a more moderate per-guest increase from a higher starting price level.

The strong coefficient on review scores (\$31.68 per point) underscores the economic value of guest satisfaction in the short-term rental market. This finding aligns with economic theory: when guests can easily compare properties and share experiences online, quality differentials become major price determinants. Hosts with consistently high ratings effectively advertise reliability and quality, justifying price premiums.

The marginal significance of review activity (negative association) warrants discussion. Higher review frequency might indicate properties with rapid guest turnover or budget-conscious pricing strategies, both of which could correlate with lower nightly rates. Alternatively, this may reflect a selection effect where expensive properties receive fewer (but highly satisfied) reviews. The small magnitude and weak significance suggest this factor plays a secondary role compared to review scores themselves.

\subsection{Addressing the Research Question}

Our initial research question sought to understand what factors optimize Airbnb pricing in Columbus. Our model provides a clear answer: competitive pricing should balance capacity utilization and room type positioning. For entire-home hosts, capacity expansion and maintaining high guest satisfaction (4.9+ stars) yield substantial price premiums. For private room hosts, the pricing leverage comes primarily through review quality rather than attempting to scale capacity. Property improvement decisions should prioritize amenities and service quality, as these drive the largest dollar returns through review score premiums.

Policymakers seeking to understand short-term rental market dynamics should note that price segregation by room type is pronounced and economically rational. The market appears to function efficiently in allocating different property types to different guest segments, with pricing reflecting genuine differences in utility and guest preferences.

\subsection{Limitations and Future Work}

Several limitations should be acknowledged. First, our model does not incorporate spatial variation; location amenities, neighborhood walkability, and proximity to attractions significantly influence prices but are not captured in our current analysis. Future work integrating geographic information systems (GIS) data or explicit location fixed effects could substantially improve model fit. Second, temporal dynamics are absent; seasonal pricing variations and platform-level changes over time were not examined. Third, amenities, while present in the raw data, were not systematically encoded into binary features; expanding this aspect could uncover additional pricing drivers.

The heteroscedasticity identified in our diagnostics suggests that advanced modeling techniques such as quantile regression might provide richer insights into pricing patterns at different market segments. Additionally, the modest $R^2$ (0.48) indicates that host-level factors (reputation, response time, cancellation policies) not available in this cleaned dataset likely play important roles in pricing decisions.

\subsection{Conclusion}

This analysis demonstrates that Airbnb pricing in Columbus is primarily determined by property characteristics (capacity and bedroom count), guest satisfaction, and room type. The significant room-type--capacity interaction reveals that the market operates efficiently with distinct pricing mechanisms for different property types. The strong effect of review scores highlights the importance of quality and reliability in generating price premiums. These findings provide guidance for hosts optimizing their listings and offer a foundation for further research into short-term rental market dynamics in urban centers.

%----------------------------------------------------------

\end{document}
